\section{Introduction}

Most medical studies rely on structured data, which often lacks
important clinical details~\cite{kho2011,Seinen2025}. For Estonian
data, free-text clinical notes have enabled the extraction of adverse
drug events for pharmacogenetic research~\cite{sarg2025} and detailed
descriptions of drug discontinuation
events~\cite{suvalov-2026-extract}. Both studies used large language
models (LLMs) for limited-scope extraction of information related to
specific medication groups. In contrast, we investigate whether LLMs
can be applied to the entire free-text EHR corpus without relying on
structured data to preselect or filter relevant records. Our
experiments are conducted on a large-scale real-world clinical
dataset, Est-Health~30~\cite{reisberg2026eh30}.

In this study, we extend the work of \v{S}uvalov \emph{et
al.}~\cite{suvalov-2026-extract} by creating a heterogeneous dataset
of drug discontinuation events and by assessing the number of
annotations required to train different LLM-based methods. As manual
prompt engineering with carefully selected few-shot examples is often
brittle and sensitive to wording~\cite{brown2020language,liu2024gpt},
we try alternatives. Fine-tuning can improve performance, but may
suffer from catastrophic forgetting when only a few hundred
annotations are available. Also, the resulting models are often
gigabytes in size. Prompt tuning offers a lightweight alternative and
has been shown to improve performance in multilingual transfer on
low-resource languages
\cite{philippy-etal-2024-soft,lester2021power}. Moreover, the
resulting soft prompts are interpretable and only a few kilobytes in
size, making them easy to share and publish with minimal privacy risk.
